% =============================================================================
%  Guarded Bulk File Reorganization
%  Ian Too
%
%  Compile with any of:
%     pdflatex tidy-up.tex   (twice, for cross-references)
%     lualatex tidy-up.tex   (recommended: full Unicode, including CJK samples)
%     xelatex  tidy-up.tex
%
%  The preamble detects the engine and configures Unicode handling to match.
% =============================================================================
\documentclass[11pt,a4paper]{article}

% ------------------------------------------------------------------ encoding
\usepackage{iftex}
% TeX counts \else and \fi even inside the branch it is skipping, so neither
% branch below may contain a nested conditional. Font selection is therefore
% recorded here and acted on after the conditional closes.
\ifPDFTeX
  \usepackage[utf8]{inputenc}
  \usepackage[T1]{fontenc}
  \usepackage{lmodern}
  \usepackage{textcomp}
  % pdfTeX cannot typeset arbitrary scripts from a Latin font. Samples outside
  % Latin are shown transliterated, with a marker, so the point survives even
  % when the glyphs cannot.
  \newcommand{\unisample}[2]{\texttt{#2}\,\textsuperscript{\tiny(non-Latin)}}
\else
  \usepackage{fontspec}
  \defaultfontfeatures{Ligatures=TeX}
  \newcommand{\unisample}[2]{\texttt{#2}\,\textsuperscript{\tiny(non-Latin)}}
  % Latin Modern has no CJK coverage. Try the usual faces on Linux, macOS and
  % Windows in turn; \providecommand means the first one found wins.
  \newcommand{\trycjkfont}[1]{\IfFontExistsTF{#1}{\providecommand{\cjkfontname}{#1}}{}}
  \trycjkfont{Noto Serif CJK JP}
  \trycjkfont{Noto Sans CJK JP}
  \trycjkfont{Hiragino Mincho ProN}
  \trycjkfont{Yu Mincho}
  \trycjkfont{MS Mincho}
  \trycjkfont{SimSun}
\fi

% Only reached with a Unicode engine that found a CJK face; otherwise the
% transliterating definition above stands and the document still builds.
\makeatletter
\@ifundefined{cjkfontname}{}{%
  \newfontfamily\cjkfont{\cjkfontname}%
  \renewcommand{\unisample}[2]{{\cjkfont #1}}%
}
\makeatother

\usepackage[english]{babel}
\usepackage{csquotes}

% -------------------------------------------------------------------- layout
\usepackage[margin=1.05in,headheight=14pt]{geometry}
\usepackage{microtype}
\usepackage{booktabs}
\usepackage{tabularx}
\usepackage{multirow}
\usepackage{array}
\usepackage{amsmath,amssymb}
\usepackage{graphicx}
\usepackage{xcolor}
\usepackage{fancyhdr}
\usepackage{titlesec}
\usepackage{caption}
\usepackage{enumitem}
\usepackage{listings}
\usepackage[ruled,vlined,linesnumbered]{algorithm2e}
\usepackage[colorlinks=true,allcolors=darkblue,breaklinks=true]{hyperref}
\usepackage{cleveref}

\definecolor{darkblue}{RGB}{15,60,125}
\definecolor{codebg}{RGB}{248,248,246}
\definecolor{codekw}{RGB}{140,30,110}
\definecolor{codecm}{RGB}{95,110,95}
\definecolor{codest}{RGB}{30,110,60}
\definecolor{rulegrey}{RGB}{120,120,120}

% ---------------------------------------------------------- arXiv-style head
\captionsetup{font=small,labelfont=bf}
\titleformat{\section}{\normalfont\large\bfseries}{\thesection}{0.6em}{}
\titleformat{\subsection}{\normalfont\normalsize\bfseries}{\thesubsection}{0.6em}{}
\titleformat{\subsubsection}{\normalfont\normalsize\itshape}{\thesubsubsection}{0.6em}{}

\pagestyle{fancy}
\fancyhf{}
\renewcommand{\headrulewidth}{0.4pt}
\fancyhead[L]{\footnotesize\scshape Guarded Bulk File Reorganization}
\fancyhead[R]{\footnotesize\thepage}

% ORCID mark: use the proper logo when the package is installed, a plain link
% otherwise, so the document builds on a bare TeX distribution.
\newif\ifhasorcidlink
\IfFileExists{orcidlink.sty}{\hasorcidlinktrue}{\hasorcidlinkfalse}
\ifhasorcidlink
  \usepackage{orcidlink}
  \newcommand{\orcidmark}[1]{\,\orcidlink{#1}}
\else
  \newcommand{\orcidmark}[1]{\,{\footnotesize\href{https://orcid.org/#1}{[ORCID]}}}
\fi

% ------------------------------------------------------------------ listings
\lstdefinelanguage{RustLang}{
  keywords={fn,let,mut,pub,struct,enum,impl,match,if,else,for,while,loop,return,
            use,mod,const,static,trait,where,self,Self,as,in,ref,move,dyn,crate,
            true,false,Some,None,Ok,Err,continue,break},
  keywordstyle=\color{codekw}\bfseries,
  ndkeywords={u32,u64,i64,usize,bool,str,String,PathBuf,Path,Vec,Option,Result,
              BTreeMap,BTreeSet,HashSet},
  ndkeywordstyle=\color{darkblue},
  sensitive=true,
  comment=[l]{//},
  morecomment=[s]{/*}{*/},
  commentstyle=\color{codecm}\itshape,
  stringstyle=\color{codest},
  morestring=[b]",
}
\lstset{
  language=RustLang,
  basicstyle=\ttfamily\footnotesize,
  backgroundcolor=\color{codebg},
  frame=leftline,
  framerule=1.1pt,
  rulecolor=\color{rulegrey},
  framesep=7pt,
  xleftmargin=10pt,
  breaklines=true,
  breakatwhitespace=true,
  showstringspaces=false,
  columns=fullflexible,
  captionpos=b,
  aboveskip=1.0em,
  belowskip=1.0em,
}
\lstdefinestyle{shell}{
  language={},
  basicstyle=\ttfamily\footnotesize,
  keywordstyle={},
  commentstyle={},
}

\newcommand{\code}[1]{\texttt{\small #1}}
\newcommand{\tool}{\textsc{tidy-up}}

% ------------------------------------------------------------------ metadata
\hypersetup{
  pdftitle={Guarded Bulk File Reorganization: System-Directory Detection, Survivable Execution, Reversible Re-grouping, and Safe Delegation to Autonomous Agents},
  pdfauthor={Ian Too},
  pdfsubject={Systems software; file management; software safety},
  pdfkeywords={file system, safety, cross-platform, Rust, undo, observability, agents, capability security, Model Context Protocol},
}

\begin{document}

% =============================================================================
\begin{center}
  {\LARGE\bfseries Guarded Bulk File Reorganization:\\[0.25em]
   System-Directory Detection, Survivable Execution, Reversible\\[0.25em]
   Re-grouping, and Safe Delegation to Autonomous Agents\par}
  \vspace{1.4em}
  {\large Ian Too\orcidmark{0009-0000-4888-1941}\par}
  \vspace{0.4em}
  {\small\href{https://orcid.org/0009-0000-4888-1941}{\texttt{https://orcid.org/0009-0000-4888-1941}}\par}
  \vspace{0.2em}
  {\small Independent Researcher\\
   \href{https://github.com/ianktoo/tidy-up}{\texttt{github.com/ianktoo/tidy-up}}\par}
  \vspace{1.6em}
\end{center}

\begin{center}
\begin{minipage}{0.86\textwidth}
\rule{\textwidth}{0.8pt}\\[0.5em]
{\centering\textbf{Abstract}\par}
\vspace{0.4em}
\small
Tools that reorganize files in bulk are, by construction, capable of destroying an
operating system installation: the same loop that sorts a downloads folder into
categories will happily scatter \code{C:\textbackslash Windows} across a dozen
subdirectories. Existing utilities address this weakly or not at all, relying on the
user to type the right path. We present the design, implementation and evaluation of
a guarded bulk reorganization system built into \tool{}, a cross-platform
command-line tool written in Rust. Four contributions are described. First, a
\emph{pure, platform-parameterized classifier} that judges whether a directory is
managed by the operating system, returning a three-level risk verdict with
machine-readable reasons; because the target platform is an argument rather than a
compile-time constant, all three rule tables (Windows, macOS, Linux) are exercised
from a single host, and a 24-entry negative corpus guards against warning fatigue.
Second, an \emph{empirical writability probe} that sidesteps the well-known
disagreement between POSIX mode bits, access-control lists, read-only mounts and
macOS System Integrity Protection by attempting a write and observing the result.
Third, a \emph{survival architecture} in which per-item failures are classified,
collected and reported rather than propagated, with a single stated exception: a
change that cannot be journaled is fatal, because a change that cannot be journaled
cannot be undone. Fourth, an \emph{idempotent re-grouping algorithm} that unpacks an
arbitrarily nested directory tree and re-files it under an ordered list of grouping
keys, computes which directories the operation will empty, and records their removal
so that undo restores the original tree exactly, including directories that held
nothing. Fifth, an account of \emph{delegating filesystem mutation to an
autonomous agent}, which we argue is admissible only where every action is
reviewable before it happens and reversible after. We contrast two interface
shapes for such delegation, validating the operations an agent names against
never letting it name one, and report that the second is strictly stronger,
because validation must anticipate every way a path can escape a root while
the alternative admits no path at all. Applying that to a Model Context
Protocol server yields an interface in which indirect prompt injection has no
argument to express itself through. We further report a structural limit of
journal-based undo: a journal records transitions, so it can answer neither
\enquote{what was left untouched} nor \enquote{can this still be undone},
both of which a reviewer needs.
The implementation adds 10{,}265 lines across sixteen modules and introduces
\emph{zero} new dependencies, keeping the transitive dependency count at 54 crates.
An evaluation over 455 automated tests reports the properties each mechanism is
designed to guarantee, including a byte-for-byte reorganize/undo round trip, the
non-overridability of a genuine permission wall, and the refusal of a plan that
reaches outside its own root.
\\[0.6em]
\textbf{Keywords:} file system safety; bulk reorganization; cross-platform systems;
capability security; agent tool design; prompt injection; idempotence;
reversible operations; observability; Rust.
\\[0.3em]
\rule{\textwidth}{0.8pt}
\end{minipage}
\end{center}

\vspace{1.2em}

% =============================================================================
\section{Introduction}
\label{sec:intro}

A bulk file organizer reads a directory, decides where each file belongs, and moves
it there. The mechanism is simple and the failure mode is severe. Nothing in the
mechanism distinguishes a downloads folder from a system directory: both are
directories full of files, and the loop that tidies one will tidy the other. A user
who types the wrong path, a script that interpolates an empty variable into
\code{organize "\$DIR"}, or an administrator experimenting with elevated privileges
can, in a single command, relocate the contents of an operating system installation
into category subdirectories.

The conventional defences are inadequate in a specific and instructive way. Dry-run
modes, confirmation prompts, refusal to overwrite, and undo journals all protect the
user against mistakes made \emph{inside} a correctly chosen directory. None of them
protects against choosing the wrong directory in the first place. An undo journal in
particular offers false comfort: replaying it restores file locations, but it cannot
restore a system that failed to boot in the interim, and it cannot help at all if the
files that moved were the ones the undo procedure itself depends on.

This paper describes what we added to \tool{}, an existing open-source cross-platform
organizer, to close that gap, together with three further mechanisms that the gap
analysis exposed as prerequisites. \Cref{sec:related} situates the work.
\Cref{sec:overview} gives the system architecture. \Cref{sec:detect} presents the
detection classifier, which is the principal contribution. \Cref{sec:perm} treats
permission determination, \Cref{sec:survival} the survival architecture,
\Cref{sec:regroup} the re-grouping algorithm, and \Cref{sec:obs} observability.
\Cref{sec:eval} evaluates the result and \Cref{sec:discussion} discusses its limits.

\subsection{Problem statement}

Let $D$ be the set of directories on a machine. A bulk reorganizer applies a
transformation $T : D \rightarrow D$ that relocates the regular files beneath a root
$r \in D$. We require a predicate
\[
  \mathrm{risk} : D \rightarrow \{\textsf{Safe}, \textsf{Caution}, \textsf{Dangerous}\}
\]
and a policy that maps a risk level, an access intent, and a set of user-supplied
flags to one of \{proceed silently, warn, warn and confirm, refuse\}. Four properties
are desirable, and they are in tension:

\begin{description}[leftmargin=1.6em,itemsep=0.25em]
  \item[P1 (Soundness, informal).] Directories the operating system manages are not
    classified \textsf{Safe}.
  \item[P2 (Precision).] Ordinary user directories are classified \textsf{Safe}. A
    system that warns about \code{D:\textbackslash Downloads} trains users to bypass
    it, destroying P1 in practice.
  \item[P3 (Testability).] The rules for each platform can be validated without
    access to that platform and without touching real system directories.
  \item[P4 (Escapability).] A user who genuinely intends to operate on a flagged
    directory can, except where the operating system itself refuses.
\end{description}

P2 deserves emphasis. A guard is a human-factors artifact as much as a technical one.
Warning fatigue is the dominant failure mode of safety mechanisms in interactive
tools, and a false positive is therefore not a cosmetic defect but a direct attack on
the mechanism's purpose. This observation drove two design decisions reported in
\Cref{sec:detect:weak} that deliberately reduce coverage.

\subsection{Contributions}

\begin{enumerate}[leftmargin=1.6em,itemsep=0.25em]
  \item A pure, platform-parameterized system-directory classifier with
    machine-readable reasons and per-reason severity, validated by a table-driven
    corpus that exercises all three platforms from one host (\Cref{sec:detect}).
  \item An empirical writability probe that is correct in the presence of ACLs,
    read-only mounts and macOS System Integrity Protection, and a corresponding
    policy rule that a permission wall is never overridable (\Cref{sec:perm}).
  \item A survival architecture with a classified failure taxonomy, a single
    documented fatal exception, and grouped, actionable end-of-run reporting
    (\Cref{sec:survival}).
  \item An idempotent, reversible re-grouping algorithm, including the computation of
    directories emptied by a run and their restoration on undo (\Cref{sec:regroup}).
  \item An account of delegating filesystem mutation to an autonomous agent: a
    machine interface whose exit codes separate a refusal from a failure, the
    argument that an interface in which an agent cannot name a path is
    strictly stronger than one that validates the paths it names, and a
    configuration model in which discovered settings may lower the bar but
    never raise it (\Cref{sec:agent}).
  \item The observation that journal-based undo cannot, structurally, answer
    the two questions a reviewer most needs, and what has to be added to it
    (\Cref{sec:agent:review}).
  \item An evaluation, and an account of the engineering constraint --- zero new
    dependencies --- that shaped every one of these designs (\Cref{sec:eval}).
\end{enumerate}

% =============================================================================
\section{Background and related work}
\label{sec:related}

\subsection{Guarded destructive operations}

The idea of refusing a dangerous argument is old and its canonical instance is
instructive. GNU \code{rm} implements \code{--preserve-root}, enabled by default,
which refuses to operate recursively on \code{/}. The design is notable for what it
does \emph{not} attempt: it protects exactly one path, by exact comparison, and makes
no judgement about \code{/usr} or \code{/home}. This is a deliberate trade toward
precision (P2) at the expense of soundness (P1), and the reasoning is sound for a
general-purpose tool invoked by scripts thousands of times a day.

A bulk reorganizer occupies a different point on that curve. It is invoked
interactively, rarely, on a directory the user names explicitly, and its blast radius
is an entire subtree. The cost of a false positive is one extra keystroke; the cost of
a false negative is an unbootable machine. We therefore extend coverage well beyond
the filesystem root while treating every extension as a precision risk to be tested
against a negative corpus.

Package managers offer a second reference point. Both \code{dpkg} and \code{rpm}
maintain a database of files owned by the system and refuse to let an installation
overwrite them. That approach is sound and precise, but it presumes a curated
database, which does not exist on Windows or macOS and which a third-party tool cannot
construct. Our classifier approximates that database with a static rule table plus
runtime facts.

\subsection{Capability determination}

Determining whether a process may write to a directory is a recurring difficulty. The
POSIX \code{access(2)} interface answers with respect to the \emph{real} user ID,
which is usually not the question being asked; \code{faccessat} with
\code{AT\_EACCESS} corrects that, but neither call is authoritative. The manual page
for \code{access(2)} explicitly warns that its result may be invalidated between the
check and the subsequent operation, a classic time-of-check-to-time-of-use hazard,
and both calls can disagree with the filesystem on NFS, with POSIX and NFSv4
access-control lists, on read-only mounts, and under macOS System Integrity
Protection, where even a process running as \code{root} is denied. On Windows the
mode-bit model does not apply at all, and \code{FILE\_ATTRIBUTE\_READONLY} on a
directory does not mean the directory is read-only: it signals a customized folder
view and is set on entirely ordinary directories such as \code{Documents}.

The technique of resolving this by attempting the operation is folklore in systems
programming rather than a published result. We adopt it, scope it carefully so that it
never runs when the caller promised not to write, and report its behaviour.

\subsection{Capability security and agent tools}

The distinction drawn in \Cref{sec:agent:capability} between validating a name
and holding an unforgeable reference is the capability-security distinction,
which is old \cite{dennis1966} and was developed at length by Miller
\cite{miller2006}. The confused-deputy problem \cite{hardy1988} is the
canonical failure of the alternative: authority exercised on behalf of a caller
who should not have had it, because the request named a resource rather than
presenting a right to it.

What is comparatively new is the caller. Tool-using language models introduce a
caller that composes requests from text it was given, and that text may be
adversarial without the operator knowing: the indirect prompt injection setting
described by Greshake and colleagues \cite{greshake2023}. Our contribution here
is not the principle but a worked application of it to a destructive filesystem
tool, and the observation that a plan/execute split maintained for testability
turns out to supply exactly the unforgeable reference such an interface needs.

\subsection{Reversible file operations}

Journaling to enable undo is standard in filesystems and in version control. \tool{}
predates this work with an append-only JSON Lines journal recording every move. The
contribution here is the observation that a journal recording only \emph{moves} is
insufficient once an operation also \emph{deletes} directories: replaying moves
recreates directories that held files, by side effect, but silently loses directories
that held nothing. \Cref{sec:regroup:undo} describes the additional record type this
requires.

% =============================================================================
\section{System overview}
\label{sec:overview}

\tool{} is a Rust command-line tool, approximately 17{,}500 lines including tests,
organized into layers with a strict dependency direction. \Cref{tab:modules}
summarizes the modules, marking those introduced or substantially revised by this
work.

\begin{table}[t]
\centering
\small
\caption{Module structure. Lines include in-file unit tests.}
\label{tab:modules}
\begin{tabularx}{\textwidth}{@{}l l X r@{}}
\toprule
\textbf{Layer} & \textbf{Module} & \textbf{Responsibility} & \textbf{Lines} \\
\midrule
\multirow{3}{*}{classify} & \code{category} & file type from extension & 273 \\
 & \code{rules} & user ignore rules & 239 \\
 & \code{projects} & detect code repositories & 110 \\
\midrule
guard\,$^\dagger$ & \code{safety} & is this a system directory & 1{,}199 \\
 & \code{safety\_tests} & rule-table corpus & 1{,}227 \\
 & \code{commands::guard} & policy: warn, confirm, refuse & 330 \\
\midrule
survive\,$^\dagger$ & \code{outcome} & failure taxonomy and collection & 531 \\
\midrule
discover & \code{scan} & directory traversal & 505 \\
\midrule
\multirow{4}{*}{decide} & \code{plan} & organize plan & 247 \\
 & \code{dedupe} & duplicate detection & 533 \\
 & \code{compare} & multi-directory overlap & 460 \\
 & \code{regroup}\,$^\dagger$ & re-grouping plan & 387 \\
\midrule
act & \code{executor} & perform a plan & 250 \\
 & \code{fsops} & non-destructive primitives & 416 \\
\midrule
remember & \code{journal} & append-only undo log & 461 \\
 & \code{restore} & reverse a journal & 315 \\
\midrule
observe\,$^\dagger$ & \code{obs} & JSON Lines run record & 608 \\
\midrule
report\,$^\ddagger$ & \code{api} & JSON envelope, exit codes, plan files & 1{,}016 \\
\midrule
delegate\,$^\ddagger$ & \code{mcp} & Model Context Protocol server & 1{,}085 \\
\midrule
permit\,$^\ddagger$ & \code{config} & policy ceiling and default floor & 807 \\
\midrule
review\,$^\ddagger$ & \code{commands::show} & what one run did, and if it can be undone & 476 \\
\midrule
present\,$^\dagger$ & \code{commands::session} & menu-settable equivalents of the flags & 371 \\
\midrule
present & \code{cli}, \code{commands}, \code{ui} & arguments, flows, terminal & 1{,}700 \\
\bottomrule
\multicolumn{4}{@{}l@{}}{\footnotesize $^\dagger$ introduced for 0.2.0 \quad $^\ddagger$ introduced for 0.3.0 (\Cref{sec:agent})}
\end{tabularx}
\end{table}

An invariant governs the layering: \emph{detection is pure, policy is impure.} The
\code{safety} module answers the question \enquote{what kind of directory is this}
without reading the environment, consulting a clock, or writing anything; the
\code{commands::guard} module decides what to do about the answer, which requires a
terminal and a user. This separation is what makes P3 (testability) achievable, and
it recurs throughout: \code{plan}, \code{regroup} and \code{dedupe} compute
reviewable data structures that \code{executor} later carries out.

% =============================================================================
\section{System-directory detection}
\label{sec:detect}

\subsection{Interface}

The classifier is a total function with no error case:
\begin{lstlisting}
pub fn classify(path: &Path, env: &Environment, facts: &Facts) -> Assessment;
\end{lstlisting}

Returning \code{Result} was considered and rejected. A classifier that can fail is one
whose callers are tempted to write \code{unwrap\_or(Safe)}, converting every transient
error into a silent authorization. Instead, every input that could fail is hoisted
into \code{Facts}, whose fields are all \code{Option}:

\begin{lstlisting}
pub struct Facts {
    pub system_attribute: Option<bool>,  // Windows FILE_ATTRIBUTE_SYSTEM
    pub reparse_point:    Option<bool>,  // Windows junction or volume mount
    pub mount_point:      Option<bool>,  // Unix: st_dev differs from parent
    pub owned_by_me:      Option<bool>,  // Unix ownership
    pub writable:         Option<bool>,  // empirical probe, Section 5
    pub fs_type:          Option<String> // Linux /proc/self/mountinfo
}
\end{lstlisting}

A failed syscall yields \code{None}, which the classifier treats as \emph{unknown}
rather than as evidence in either direction. The corresponding test asserts that a
\code{Facts::default()} --- everything unknown --- never produces a finding on its own.

\subsection{Platform as a parameter}

The single decision that most improves testability is making the platform an
argument:

\begin{lstlisting}
pub enum Platform { Windows, MacOs, Linux, Other }
\end{lstlisting}

This forced a consequence that is easy to miss. Rust's \code{std::path} parses paths
for the \emph{build} target: on Linux, \code{Path::new("C:\textbackslash\textbackslash Windows").components()}
yields a single component, because the drive prefix is not recognized. Any classifier
built on \code{std::path::Components} is therefore untestable across platforms by
construction. We implement a platform-parameterized splitter:

\begin{lstlisting}
fn parts(path: &Path, platform: Platform) -> Parts;   // (prefix, components)
\end{lstlisting}

\code{Parts} normalizes four things that are individually minor and collectively the
source of most real bugs in path matching:

\begin{enumerate}[leftmargin=1.6em,itemsep=0.2em]
  \item \textbf{Verbatim prefixes.} \code{\textbackslash\textbackslash?\textbackslash C:\textbackslash Windows} and
    \code{C:\textbackslash Windows} denote the same directory, as do
    \code{\textbackslash\textbackslash?\textbackslash UNC\textbackslash srv\textbackslash sh} and
    \code{\textbackslash\textbackslash srv\textbackslash sh}. An upstream function in \tool{} deliberately
    preserves the UNC verbatim form, so the classifier must fold both spellings.
  \item \textbf{Case.} Windows and macOS compare case-insensitively by default; Linux
    does not. Components are case-folded only on the former.
  \item \textbf{Trailing punctuation.} Windows silently strips trailing dots and
    spaces from names, so \code{C:\textbackslash Windows.\textvisiblespace} resolves to
    \code{C:\textbackslash Windows} and must match the same rule.
  \item \textbf{Dot segments.} \code{.} is dropped and \code{..} pops, so a path is
    judged by where it points.
\end{enumerate}

Matching is then \emph{component-wise}, never textual. The distinction is not
pedantic: \code{C:\textbackslash Program Files Custom} has \code{C:\textbackslash Program Files} as a string
prefix but is an entirely unrelated directory, and \code{/usrlocal} is not inside
\code{/usr}. Both cases are asserted by name in the test suite.

Note also that upstream canonicalization closes an evasion path: because the target is
resolved before classification, a symbolic link or NTFS junction pointing at
\code{/usr} is judged as \code{/usr}. The original spelling is retained only for the
message shown to the user.

\subsection{Rule structure}

Each platform has a table of entries and, for the Unix platforms, a list of
\emph{carve-outs} evaluated first. An entry is:

\begin{lstlisting}
struct Entry {
    comps:    &'static [&'static str],  // components below the root
    exact:    bool,                     // this directory only, or the subtree
    severity: Risk,
    label:    &'static str,             // human description
}
\end{lstlisting}

A carve-out either declares a subtree \textsf{Safe} with respect to the path rules or
downgrades it to \textsf{Caution}. \Cref{tab:rules} gives the substantive content.

\begin{table}[t]
\centering
\small
\caption{Rule tables. \textbf{D}~=~\textsf{Dangerous} (refused), \textbf{C}~=~\textsf{Caution}
(warn and confirm), \textbf{S}~=~carve-out suppressing broader rules. Carve-outs are
evaluated first.}
\label{tab:rules}
\begin{tabularx}{\textwidth}{@{}l l X@{}}
\toprule
\textbf{Platform} & \textbf{Level} & \textbf{Locations} \\
\midrule
\multirow{4}{*}{Windows}
 & D & \code{\%SystemRoot\%} (fallbacks \code{\%windir\%},
       \code{<drive>\textbackslash Windows}); \code{\%ProgramFiles\%},
       \code{(x86)}, \code{\%ProgramW6432\%}; \code{\%ProgramData\%};
       \code{<drive>\textbackslash Users} (the container); \code{\%USERPROFILE\%} itself;
       \code{AppData}; \code{\$Recycle.Bin}; \code{System Volume Information};
       \code{Recovery}; \code{Config.Msi}; \code{PerfLogs}; \code{\$WinREAgent};
       drive roots; UNC share roots \\
 & C & \code{\%PUBLIC\%}; a bare reparse point \\
 & S & anything strictly inside \code{\%TEMP\%} \\
\midrule
\multirow{3}{*}{macOS}
 & D & \code{/}; \code{/System} (incl. \code{/System/Volumes/*}); \code{/Library};
       \code{/Applications}; \code{/usr}; \code{/bin}; \code{/sbin}; \code{/etc};
       \code{/var}; \code{/private}; \code{/cores}; \code{/Network}; \code{/.vol};
       \code{/Volumes} and \code{/Users} as containers; \code{\$HOME} itself;
       \code{\$HOME/Library} \\
 & C & \code{/opt}; \code{/usr/local/*}; \code{\$HOME/Applications};
       \code{\$HOME/.Trash}; a bare mount point \\
 & S & \code{/Volumes/<disk>/*}; \code{/Users/<user>/*};
       \code{/System/Volumes/Data/Users/<user>/*}; \code{/private/tmp}; \code{/tmp};
       \code{\%TEMP\%} \\
\midrule
\multirow{3}{*}{Linux}
 & D & \code{/}; \code{/bin}; \code{/sbin}; \code{/lib*}; \code{/usr}; \code{/etc};
       \code{/boot}; \code{/dev}; \code{/proc}; \code{/sys}; \code{/run};
       \code{/snap}; \code{/nix}; \code{/efi}; \code{/lost+found}; \code{/var};
       \code{/home} (container); \code{/root}; \code{\$HOME} itself; any
       pseudo-filesystem \\
 & C & \code{/opt}; \code{/srv}; \code{/usr/local/*}; \code{/media} and \code{/mnt}
       as containers \\
 & S & \code{/tmp}; \code{/var/tmp}; \code{/media/<u>/<v>/*}; \code{/mnt/<x>/*};
       \code{/run/media/<u>/<v>/*} \\
\bottomrule
\end{tabularx}
\end{table}

\subsubsection{Windows: variables augment, never replace}

Environment variables are read once, but the name-based rules are applied relative to
\emph{any} drive root regardless. A scrubbed environment, a service account, or a
second Windows installation on \code{D:} therefore remains protected, while
\code{\%SystemRoot\%} pointing at an unconventional location adds coverage rather than
moving it. Two tests state this directly: one clears the variable map entirely and
asserts the name rules still fire; the other sets \code{SystemRoot=E:\textbackslash Win} and asserts
that both \code{E:\textbackslash Win\textbackslash System32} and \code{C:\textbackslash Windows} are refused.

\subsubsection{macOS: firmlinks and the data volume}

Two macOS-specific hazards required carve-outs. \code{/usr/local} is writable and
full of user-installed software, so the \code{/usr} rule is downgraded there rather
than applied. More subtly, on APFS volumes a user's home directory is reachable both
as \code{/Users/me} and as \code{/System/Volumes/Data/Users/me}; canonicalization
usually yields the short form, but not under all mount configurations. Without the
second carve-out, a user's own \code{Downloads} directory would be refused as part of
\code{/System}.

System Integrity Protection has no public query interface. The
\code{rootless\_check\_trusted} symbol is private SPI in \code{libSystem} and linking
against it would be both unstable and inappropriate for a third-party tool. We treat
the path table as a practical superset and rely on the writability probe
(\Cref{sec:perm}) as the backstop: SIP-protected paths reject writes even for
\code{root}, so \code{NotWritable} fires, and that reason is deliberately
non-overridable.

\subsubsection{Linux: pseudo-filesystems}

Rather than enumerate kernel-virtual mount points, we determine the filesystem type
by parsing \code{/proc/self/mountinfo}. The parser is a pure function of the file
contents and the target path, which permits fixture-based testing on any host:

\begin{lstlisting}
pub(crate) fn fs_type_from_mountinfo(text: &str, path: &Path) -> Option<String>;
\end{lstlisting}

The format requires care. Field 5 is the mount point, octal-escaped so that
\code{\textbackslash 040} denotes a space; a variable number of optional fields follows; a literal
\code{-} separates them from the filesystem type. The longest mount point that is a
component-wise prefix of the target wins, which resolves nested mounts correctly.
Twenty-one pseudo-filesystem types are refused.

\code{tmpfs} is deliberately \emph{absent} from that set. \code{/tmp} and
\code{/dev/shm} are \code{tmpfs}, and \code{/tmp} is a perfectly reasonable directory
to tidy. Including it would have been the more \enquote{complete} choice and the
wrong one.

\subsection{Composition}

\Cref{alg:classify} gives the composition. Three details are load-bearing.

\begin{algorithm}[t]
\SetAlgoLined
\DontPrintSemicolon
\KwIn{path $p$, environment $e$, facts $f$}
\KwOut{assessment with risk level and ordered reasons}
$P \leftarrow \mathrm{parts}(p, e.\mathrm{platform})$\;
$R \leftarrow \emptyset$\;
$\mathit{inTemp} \leftarrow P$ strictly inside $\mathrm{parts}(e.\mathrm{temp})$\;
\If{$P$ has a prefix and no components}{
  $R \leftarrow R \cup \{\textsf{UncShareRoot}$ if UNC else $\textsf{FilesystemRoot}\}$\;
}
\If{$P = \mathrm{parts}(e.\mathrm{home})$}{$R \leftarrow R \cup \{\textsf{HomeRoot}\}$\;}
\lIf{$\neg\,\mathit{inTemp}$}{$R \leftarrow R \cup \mathrm{pathRules}(P, e)$}
\If{$f.\mathrm{fs\_type} \in \mathit{PSEUDO}$}{$R \leftarrow R \cup \{\textsf{PseudoFilesystem}\}$\;}
\lIf{$f.\mathrm{system\_attribute} = \textsf{Some}(\top)$}{$R \leftarrow R \cup \{\textsf{SystemAttribute}\}$}
\lIf{$f.\mathrm{writable} = \textsf{Some}(\bot)$}{$R \leftarrow R \cup \{\textsf{NotWritable}\}$}
\lIf{$f.\mathrm{owned\_by\_me} = \textsf{Some}(\bot)$}{$R \leftarrow R \cup \{\textsf{NotOwned}\}$}
\If{$(f.\mathrm{mount\_point} \vee f.\mathrm{reparse\_point}) \wedge R = \emptyset$}{
  $R \leftarrow \{\textsf{MountPoint}\}$ \tcp*{only when nothing stronger fired}
}
sort $R$ by severity, descending\;
remove $r \in R$ whose explanation duplicates an earlier one\;
\Return $(\max_{r \in R} \mathrm{severity}(r), R)$, or \textsf{Safe} if $R = \emptyset$\;
\caption{\textsc{Classify}}
\label{alg:classify}
\end{algorithm}

\paragraph{The temporary-directory carve-out is load-bearing.}
On two of three platforms the system temporary directory lies inside a protected
subtree: Windows places it under \code{\%LOCALAPPDATA\%\textbackslash Temp}, and macOS
canonicalizes it to \code{/private/var/folders/...}, matching both \code{/private} and
\code{/var}. Without the carve-out the guard would refuse the system temporary
directory --- a legitimate target --- and, because the project's entire pre-existing
test suite builds its fixtures with \code{tempfile}, would have broken every test on
first commit. The dependency between a safety rule and the test infrastructure was not
anticipated during design and was discovered only when enumerating the rule table.

\paragraph{Mount points are suppressed when anything stronger fires.}
Every external drive is a mount point. Reporting it alongside \enquote{this is
\code{/usr}} adds noise without information, so the reason is emitted only when it is
the sole finding, and then only at \textsf{Caution}.

\paragraph{Reasons are deduplicated by rendered explanation.}
\code{C:\textbackslash Windows} matches the \code{\%SystemRoot\%} rule, the \code{\%windir\%} rule and
the name rule simultaneously. The first implementation printed three near-identical
lines. Deduplicating structurally is insufficient, because the three findings differ
in their \code{matched} field; deduplicating by the sentence the user actually reads
is both simpler and correct. This was found by running the tool against the real
\code{C:\textbackslash Windows}, not by any test, and a regression test was added afterwards.

\subsection{Deliberately weakened rules}
\label{sec:detect:weak}

Two candidate rules were weakened or removed on precision grounds (P2).

\emph{\enquote{Outside the home directory}} was removed entirely. It is a plausible
heuristic and would have been catastrophic in practice:
\code{D:\textbackslash Downloads}, \code{/mnt/photos} and \code{/Volumes/Backup} are the
\emph{normal} targets for this class of tool. A rule that fires on the normal case
teaches users to pass the override flag reflexively, after which the guard protects
nobody.

\emph{\enquote{Shallow path}} (a directory one level below a filesystem root) was
retained only in conjunction: it fires only when the directory is additionally not
under the user's home and demonstrably not owned by the current user, and even then
only at \textsf{Caution}. \code{/data}, \code{/srv2} and \code{D:\textbackslash Media} are ordinary
directories that a naive depth rule would flag.

The general principle we extracted: \emph{composition, not any single weak signal.} A
directory that is shallow \emph{and} unowned \emph{and} unwritable \emph{and} a mount
point is almost certainly system-managed; any one of those alone is not.

\subsection{Policy}
\label{sec:detect:policy}

Policy is separated from detection and lives in \code{commands::guard}.
\Cref{tab:policy} states it completely.

\begin{table}[t]
\centering
\small
\caption{Guard policy. \code{--yes} answers a confirmation; it never substitutes for
the override flag.}
\label{tab:policy}
\begin{tabularx}{\textwidth}{@{}l X X X@{}}
\toprule
\textbf{Risk} & \textbf{Reporting command} & \textbf{Destructive command} &
\textbf{With \code{--allow-system-folder}} \\
\midrule
\textsf{Safe} & silent & silent & silent \\
\textsf{Caution} & warn & warn, then confirm (default \emph{no}) & same \\
\textsf{Dangerous} & warn, proceeds & \textbf{refused}, non-zero exit, names the flag &
  warn, then confirm (default \emph{no}) \\
\textsf{NotWritable} & warn, proceeds & \textbf{refused} & \textbf{still refused} \\
\bottomrule
\end{tabularx}
\end{table}

Four aspects warrant comment.

\textbf{\code{--yes} is not an override.} This is the property the evaluation treats
as flagship. Automated invocations that pass \code{--yes} continue to work on ordinary
directories and begin failing loudly on system ones. The resulting breakage of
existing scripts is intended, and is the mechanism by which the guard reaches
unattended use at all.

\textbf{The override reaches the question, not past it.} With
\code{--allow-system-folder} the user is still asked, the default answer is
\emph{no}, and with no terminal attached the run stops. Obtaining a destructive
outcome on a flagged directory requires the flag \emph{and} an affirmative answer.

\textbf{A permission wall is not overridable.} No command-line flag grants permission
that the operating system refused, and pretending otherwise would produce a run that
fails item by item after the user has already been warned. The refusal message says
explicitly that the flag will not help.

\textbf{Reporting commands never block.} \code{analyze} and \code{history} warn and
continue. Analyzing an entire drive is an advertised use of the tool; blocking it
would have been a regression, and an existing test covers it.

Two further points of mechanism. First, the guard runs at the single choke point where
a command turns a user-supplied path into a validated root, and \emph{not} inside the
library function that performs that resolution: that function is pure, returns a
library error type, cannot prompt, and is called by multi-directory validation paths
where a policy decision would be wrong. Second, \code{distribute} guards its
destinations as well as its sources, since \code{--to C:\textbackslash Windows} is as damaging as
\code{--from}; a destination that does not yet exist is judged by the nearest existing
ancestor, which is where it would be created.

Finally, \code{--include-hidden} escalates \textsf{Caution} to \textsf{Dangerous}.
Bringing dotfiles into scope widens the blast radius into exactly the territory where
configuration and system state live; a \textsf{Safe} directory remains \textsf{Safe},
because hidden files in a user's own downloads directory are not a hazard.

\subsection{Reaching the override without a command line}
\label{sec:detect:menu}

An escape hatch that cannot be reached is not an escape hatch. \Cref{tab:policy}
promises that a determined user can proceed (P4), but the promise is carried entirely
by a command-line flag, and the tool is also usable by running it with no arguments,
which starts a menu. Someone who launched it by double-clicking has nowhere to put a
flag. For them, P4 did not hold: not only was the override unreachable, so were
preview mode, hidden-file inclusion and every ignore rule, because the menu
constructed each command with default arguments.

This is easy to miss when the guard is designed, tested and demonstrated from a
shell. The fix is a settings screen holding the same options for the duration of the
session, with whatever is not at its default printed above the menu so the next
action is never a surprise. Nothing is persisted: the tool writes no configuration
file outside the folders it processes, and that property was worth more than
remembering a preference between runs.

Two details preserve the policy rather than route around it. Enabling the override
asks for confirmation, and the guard still asks again before acting, so it remains
two deliberate steps rather than one sticky setting. And \enquote{skip
confirmations}, the menu equivalent of \code{--yes}, is deliberately absent: in an
interactive session the confirmation \emph{is} the review step, and there is no
script waiting on it.

The general point is that a safety mechanism has to be evaluated against every way
the program can be started, not only the one its author uses.

% =============================================================================
\section{Permission determination}
\label{sec:perm}

\subsection{Empirical probing}

Given the disagreements catalogued in \Cref{sec:related}, we determine writability by
attempting it:

\begin{lstlisting}
fn probe_writable(dir: &Path) -> Option<bool> {
    let name = format!(".tidy-up-write-test-{}-{nanos}-{unique}", process::id());
    match fs::File::create_new(dir.join(&name)) {
        Ok(f)  => { drop(f); let _ = fs::remove_file(dir.join(&name)); Some(true) }
        Err(e) if e.kind() == PermissionDenied
               || e.kind() == ReadOnlyFilesystem  => Some(false),
        Err(_) => None,   // out of space, directory vanished: unknown, not "no"
    }
}
\end{lstlisting}

Three properties matter. The name is unique per process, per nanosecond and per
monotonic counter, so concurrent runs cannot collide. \code{create\_new} fails rather
than truncating if the name somehow exists. Most importantly, only
\code{PermissionDenied} and \code{ReadOnlyFilesystem} produce \code{Some(false)}:
every other error yields \code{None}, because an unknown must never be reported as an
accusation. A test asserts exactly this for a directory that does not exist.

The probe is \emph{scoped}. It runs only for commands about to write, and never under
\code{--dry-run}. A reporting command and a dry run both promise to leave the disk
untouched, and that promise outranks the diagnostic value of knowing writability for
an operation that was not going to write. Two tests assert that the directory is still
empty after a read-only assessment and after a dry-run assessment.

\subsection{Ownership, and a dependency declined}

Unix ownership is available cheaply, but the standard library exposes no way to read
the current \emph{effective} user ID, so comparison requires \code{libc} or
\code{rustix}. We shipped without ownership, retaining
\code{Facts::owned\_by\_me: Option<bool>} so it can be supplied later without changing
the classifier or any test. The justification is that the writability probe already
answers the question that determines the outcome, that \code{NotOwned} is only ever
\textsf{Caution}, and that the project's zero-dependency constraint
(\Cref{sec:eval:deps}) is a stated property of the artifact. Should ownership prove
necessary, \code{rustix} is preferred over \code{libc} because it is implemented in
pure Rust on \code{linux-raw} and preserves the dependency-free static musl build.

\subsection{Reporting refusals distinctly}

A permission failure previously surfaced as \code{/srv/data: Permission denied
(os error 13)} or, below the root, as the single word \enquote{unreadable}. We
separated it throughout: a new \code{Error::Denied} variant renders as
\code{cannot read <path>: permission denied} and permits the command layer to add an
actionable hint, and \code{SkipReason} gained a \code{Denied} variant distinct from
\code{Unreadable}. The difference is the difference between a user concluding
\enquote{\tool{} is broken} and \enquote{I do not own that directory}.

\subsection{Hidden and system-owned files}

Three defects in hidden-file handling were found and corrected.

\begin{enumerate}[leftmargin=1.6em,itemsep=0.2em]
  \item \textbf{macOS \code{UF\_HIDDEN}.} Files hidden with \code{chflags hidden}
    carry a flag in \code{st\_flags} that no leading-dot check detects.
    \code{std::os::macos::fs::MetadataExt::st\_flags} is stable, so this cost nothing
    but the recognition that it was missing.
  \item \textbf{Windows \textit{hidden} conflated with \textit{system}.} The previous
    implementation tested the mask \code{0x2 | 0x4} as a unit, so
    \code{--include-hidden} also un-skipped every file Windows marks as its own.
    Asking to see dotfiles is not asking to move \code{pagefile.sys}. The attributes
    are now tested separately, a \code{SkipReason::SystemAttribute} reports the
    second, and it is \emph{not} governed by \code{--include-hidden}.
  \item \textbf{Escalation.} As noted above, including hidden items raises a
    \textsf{Caution} verdict to \textsf{Dangerous}.
\end{enumerate}

% =============================================================================
\section{Survival architecture}
\label{sec:survival}

\subsection{Principle}

A bulk operation over thousands of files will encounter a locked file, an unreadable
directory, or a name the filesystem rejects. Aborting costs the user the entire
remaining run; worse, it leaves the operation half-complete in a way the user did not
choose. The principle adopted is:

\begin{quote}
\itshape
Fail the run only when continuing would do something the user did not ask for, or
would lose the ability to undo it. Everything else is a skip: classify it, record it,
step over it, and report it at the end.
\end{quote}

\subsection{Failure taxonomy}

A single type carries every failure, and one function classifies every I/O error, so
that counting is consistent across modules:

\begin{lstlisting}
pub struct Problem { path: PathBuf, op: Op, cause: Cause, message: String }

pub enum Cause { Denied, NotFound, InUse, Exists, CrossesDevices,
                 NoSpace, Invalid, Panic, Other }
\end{lstlisting}

\code{Cause::of} maps \code{io::ErrorKind} to a cause in one place. Two mappings are
worth noting. \code{ReadOnlyFilesystem} maps to \code{Denied}, because from the user's
point of view it is a permission problem with the same remedy. Windows reports a file
held open by another process as \code{ERROR\_SHARING\_VIOLATION} (32) or
\code{ERROR\_LOCK\_VIOLATION} (33), neither of which has a stable
\code{io::ErrorKind}; these are recognized by raw code and mapped to \code{InUse}, or
every locked file would be reported as \enquote{other}.

Two combinators express the principle directly:

\begin{lstlisting}
pub fn try_or_skip<T>(&mut self, path: &Path, op: Op,
                      f: impl FnOnce() -> Result<T>) -> Option<T>;
pub fn catch_or_skip<T>(&mut self, path: &Path, op: Op,
                        f: impl FnOnce() -> Result<T> + UnwindSafe) -> Option<T>;
\end{lstlisting}

The second additionally contains panics via \code{catch\_unwind}, converting a bug in
the processing of one file into a \code{Cause::Panic} problem rather than the loss of
the entire run. A test panics inside the closure and asserts both that the problem is
recorded with its message and that the collection remains usable afterwards.

\subsection{Converting the traversal}

The directory walk previously propagated \code{Result} through its recursion, so a
failure anywhere in a large tree aborted the scan. It is now infallible below the
root: every failure becomes a \code{Skipped} entry with a classified reason. Only the
root itself can fail the command, which is correct, because a root that cannot be
listed offers nothing to do.

A test constructs a tree with one subdirectory at mode \code{0o000} and asserts both
that every other file is still returned and that the failed directory is named exactly
once with reason \code{Denied}. It self-disables when running as \code{root}, for whom
nothing is unreadable.

\subsection{The one fatal exception}

Journal write failure remains fatal, and the in-flight move is rolled back. This is
the deliberate exception to the principle: a change that cannot be recorded cannot be
undone, and silently continuing would strand files with no way back. Stating the
exception explicitly, and having exactly one, is what makes the rest of the policy
legible.

\subsection{Reporting}

Reporting is grouped by cause with an actionable hint per group, then individual
paths, capped unless \code{--verbose} is given:

\begin{lstlisting}[style=shell]
!  23 items could not be processed:
   Permission denied               18  (pick a folder you own, or run as an administrator)
   In use by another program        4  (close the programs holding these files and run again)
   No longer there                  1  (these were deleted or renamed while tidy-up was running)
   Downloads/locked.pdf: ...
   ... and 13 more (use --verbose to list them all)
\end{lstlisting}

Twenty identical lines reading \enquote{permission denied} inform nobody; one line
reading \enquote{18 permission denied, pick a folder you own} is actionable. A run
that skipped items still exits zero, because it performed the work it could;
\code{--strict} makes it exit non-zero for callers that need to know.

% =============================================================================
\section{Reversible re-grouping}
\label{sec:regroup}

\subsection{Motivation and interface}

The pre-existing \code{organize} command sorts loose files into category directories
and deliberately refuses to descend into its own output, scanning one level deep and
skipping known category names. That is correct for a downloads directory and useless
for a directory somebody organized badly years ago, which is the common real case.

\code{reorganize} inverts both choices. It descends without limit, including into
directories a previous run created, re-files every file under an ordered list of
grouping keys, and deletes the directories it empties.

\begin{lstlisting}[style=shell]
tidy-up reorganize <path> --by year,type      # 2024/Images/photo.jpg
tidy-up reorganize <path> --by type,ext       # Images/PNG/photo.png
tidy-up reorganize <path> --by alpha
\end{lstlisting}

\Cref{tab:keys} lists the keys. Each becomes one level of nesting, in the order given.

\begin{table}[t]
\centering
\small
\caption{Grouping keys. Every key is total: no file is left unplaced.}
\label{tab:keys}
\begin{tabularx}{\textwidth}{@{}l l X@{}}
\toprule
\textbf{Key} & \textbf{Example directory} & \textbf{Derivation} \\
\midrule
\code{type}  & \code{Images}, \code{3D Models} & shared with \code{organize}, so the
  two commands agree \\
\code{ext}   & \code{PDF}, \code{No extension} & upper-cased extension; the
  \enquote{document type} grouping \\
\code{year}  & \code{2024}, \code{Unknown date} & modification time \\
\code{month} & \code{03-March} & numeric prefix so directories sort in calendar order \\
\code{day}   & \code{2024-03-17} & modification time \\
\code{size}  & \code{Small (under 10 MiB)} & five bands, powers of 1024 \\
\code{alpha} & \code{A}, \code{0-9}, \code{Other} & first character, upper-cased \\
\bottomrule
\end{tabularx}
\end{table}

Totality is enforced by test: a corpus of awkward inputs --- no extension, no
timestamp, a leading digit, an accented initial (\code{étude.mp3}), a non-Latin
initial (\unisample{日本語.txt}{nihongo.txt}), and the name \code{...} --- is passed
through every key, and each is asserted to produce a non-empty, already-sanitized
name.

\subsection{Name sanitization}

Directory names derived from file contents can contain anything. Sanitization
replaces the nine characters Windows rejects and all control characters, trims
trailing dots and spaces, falls back to \code{Other} if nothing survives, and appends
an underscore to the 22 reserved device names. The last case is not hypothetical:
\code{driver.con} is an ordinary file name, and grouping by extension would otherwise
attempt to create a directory named \code{CON}, which Windows will not create at any
path.

\subsection{Idempotence}

\code{organize} obtains idempotence from its skip list. \code{reorganize} has no skip
list and must compute it: for each file, the destination is derived, and if it equals
the current path no move is recorded. A move from a path to itself would otherwise
churn the journal on every run and never converge.

The property is asserted twice, at two levels. A unit test feeds the destinations of
one plan back in as the input of a second and asserts the second plan is empty. An
integration test runs the compiled binary twice over a fixture and asserts that the
second run reports nothing to do and that the file and directory listings are
byte-identical.

\subsection{Computing emptied directories}

\Cref{alg:emptied} computes which directories the run will leave empty. The subtlety
is that emptiness is transitive upward --- a directory containing only directories that
become empty is itself empty --- while any retained item disqualifies a directory
\emph{and all of its ancestors}.

\begin{algorithm}[t]
\SetAlgoLined
\DontPrintSemicolon
\KwIn{root $r$; $V$ = parents of moved items; $K$ = ancestors of retained items}
\KwOut{directories to remove, deepest first}
$C \leftarrow \emptyset$\;
\ForEach{$d \in V$}{
  \While{$d \neq r$ \textbf{and} $d$ is under $r$}{
    $C \leftarrow C \cup \{d\}$; \quad $d \leftarrow \mathrm{parent}(d)$\;
  }
}
$C \leftarrow C \setminus K$ \tcp*{anything holding a survivor, and its ancestors}
$C \leftarrow C \setminus \{d : d \text{ is or is under } r/\texttt{.tidy-up}\}$\;
sort $C$ by component count, descending\;
\Return $C$\;
\caption{\textsc{EmptiedDirectories}}
\label{alg:emptied}
\end{algorithm}

Items are \emph{retained} when they are skipped by an ignore rule, are unreadable,
are symbolic links, or are code project directories that policy keeps in place. The
tool's own \code{.tidy-up} state directory is excluded unconditionally: it holds the
undo journals, and a run that swept away its own means of reversal would be
self-defeating.

Two safety layers back the computation. Directories are removed deepest first, and
removal uses \code{remove\_dir}, which succeeds only on a genuinely empty directory.
An error in \Cref{alg:emptied} therefore produces a directory that is not removed,
never one that is removed with contents. Consistent with \Cref{sec:survival}, a
removal that fails is recorded as a problem and the run continues.

\subsection{Undo}
\label{sec:regroup:undo}

Reversing a re-grouping requires more than replaying moves. Files whose parent
directories vanished are restored correctly by side effect, since the move primitive
creates destination parents. Directories that ended up holding \emph{nothing} have no
move to recreate them and would be silently lost.

The journal therefore gained one record type:

\begin{lstlisting}
DirRemoved { path: PathBuf }   // relative to the root, like every other record
\end{lstlisting}

written after each successful removal --- after, so that a removal which did not
happen is never recorded. Restore replays records in reverse, which places directory
recreation before the moves that repopulate the subtree.

The property is asserted end to end. An integration test constructs a deliberately
badly organized tree --- nested category directories, a \code{\_Duplicates} directory,
a code project, a dotfile, and an already-empty directory --- reorganizes it by
\code{year,type}, then restores, and asserts that both the complete file map (relative
path to contents) and the complete directory listing equal their original values. The
empty directory is in the fixture specifically because it is the case the journal
could not previously express.

\subsection{Time zones: a stated limitation}

Date grouping requires civil dates, and the Rust standard library has no time-zone
interface. Obtaining the local UTC offset portably would require a TZif parser plus
\code{/etc/localtime} on Unix and a registry read on Windows --- in practice, a
dependency.

Rather than be correct on one platform and silently wrong on another, dates are UTC,
and an environment variable \code{TIDY\_UP\_UTC\_OFFSET} accepting \code{+03:00},
\code{-0530}, \code{+7} or \code{Z} applies an offset. The consequence is real and
documented: with no offset set, a file modified at 01:00 local time in UTC+3 on
1 January is grouped under the previous year. A test states exactly this case in both
configurations. We consider a documented, correctable limitation preferable to an
undocumented platform-dependent one, and note that the effect is confined to a small
fraction of files near boundaries and is fully reversible.

% =============================================================================
\section{Observability}
\label{sec:obs}

\subsection{What the journal already provides}

An important negative design decision: the run log is \emph{not} a per-move log,
because one already exists. The undo journal records every move as it happens,
unbuffered, so that a crash leaves a usable record. Duplicating it would double the
writes to say the same thing twice.

What a journal of moves cannot record is everything that did \emph{not} happen: the
safety verdict, what was skipped and why, what failed, and how long each stage took.
That is the gap the run log fills, and the two artifacts join on a shared identifier.

\subsection{Format and policy}

Events are buffered and written once, at the end of the run, to
\code{<root>/.tidy-up/logs/<run-id>.jsonl}, beside the journal. Buffering is
acceptable precisely because the journal is the crash-safe record: a run log missing
after a crash costs a diagnosis, not data. A complete log from a real run:

\begin{lstlisting}[style=shell]
{"type":"run","command":"reorganize","version":"0.1.1","os":"windows","format":1,...}
{"type":"safety","path":"...\\logdemo","risk":"safe","reasons":[],"allowed":false}
{"type":"scan","root":"...","files":2,"bytes":3,"projects":0,"skipped":{},"ms":0}
{"type":"plan","operation":"reorganize","moves":2,"bytes":3}
{"type":"execute","journal_id":"20260928-204156","moved":2,"bytes":3,
 "dirs_removed":2,"failed":0,"ms":0}
{"type":"end","status":"ok","ms":4,"metrics":{"roots":1,"files_scanned":2,
 "bytes_scanned":3,"skipped":0,"planned_moves":2,"moved":2,"bytes_moved":3,
 "dirs_removed":2,"failed":0,"denied":0}}
\end{lstlisting}

Events are a closed Rust enum serialized with an internal tag, which makes the log a
type-checked, versioned contract rather than a collection of stringly-typed fields
that a typo can silently rename --- the same property the journal format already has.
\code{status} distinguishes \code{ok}, \code{error}, \code{cancelled} and
\code{blocked}, the last recording a guard refusal.

Logging is off by default, enabled by \code{--log} or \code{TIDY\_UP\_LOG}, preserving
the project's documented promise to write nothing the user did not ask for. Logs are
pruned to the newest twenty per directory; journals, which are undo data, are never
touched by pruning, and a test asserts this by placing a non-log file in the directory
and verifying it survives. A log that cannot be written degrades silently: a
diagnostic must never change a run's outcome or exit code, and a test asserts that a
run whose log directory path is occupied by a regular file still succeeds.

\subsection{A conceded trade-off}

The sink is process-level rather than threaded through the call graph. A
command-line tool performs exactly one run per process, and threading a handle through
nine commands to record six events would have cost more in noise and merge risk than
it returned. The cost is that tests of the sink exercise it directly or through the
compiled binary rather than in parallel against isolated instances. We consider this
the right trade at this scale and would revisit it if the library were embedded.

% =============================================================================
\section{Delegating to an autonomous agent}
\label{sec:agent}

The mechanisms so far protect a person from a mistake. This section is about a
different caller: a program that decides what to do without being told each
time. The distinction matters because an agent acts in bulk, acts quickly, and
acts on text it was given, which may not have come from the person it is
working for.

Everything here is built on a property the earlier work already had, and which
we think is the precondition for delegating filesystem mutation at all:
\emph{every action can be reviewed before it happens and reversed after}. A
\code{Plan} is pure data (\Cref{sec:overview}), and every change is journaled
(\Cref{sec:regroup:undo}). Without both, no amount of interface design makes
this safe; with both, the remaining question is what the agent is allowed to
express.

\subsection{A machine interface}
\label{sec:agent:api}

Every command accepts \code{--json} and prints exactly one object on standard
output, whatever happened:

\begin{lstlisting}[style=shell]
{"tidy_up":"0.3.0","format":1,"command":"organize","status":"ok","outcome":{...}}
\end{lstlisting}

The human and machine renderings are two views of one value, not two
implementations: a command returns an \code{Outcome} and the caller decides how
to render it. Keeping one model is what stops the two drifting.

Making that reliable required a single choke point rather than a guard at each
call site. A \code{out!} macro replaces every prose \code{println!} in the
crate, and progress bars are suppressed. Guarding each site individually would
have worked on the day and failed silently the next time somebody added a
print, which is the failure mode worth designing against.

\paragraph{Exit codes distinguish a refusal from a failure.}
Previously every non-success exited \code{1}. A caller cannot act on that: a
refusal is an answer and must not be retried, whereas a transient failure might
reasonably be. \Cref{tab:exits} gives the mapping. To make it honest, refusals
are \emph{typed} rather than worded: a \code{Refused} error carries an
\code{ErrorCode}, and both the exit code and the JSON error derive from it. An
earlier version decided the run-log status by testing whether the message began
with \enquote{refusing}, which is the kind of thing that survives until somebody
rewords a string.

\begin{table}[t]
\centering
\small
\caption{Exit codes. \code{3} is the one that did not exist before.}
\label{tab:exits}
\begin{tabularx}{\textwidth}{@{}l X@{}}
\toprule
\textbf{Code} & \textbf{Meaning} \\
\midrule
\code{0} & Success, including \enquote{nothing to do}, a cancelled confirmation, and a run that skipped items it could not process \\
\code{1} & Something went wrong and stopped the run \\
\code{2} & The arguments did not make sense \\
\code{3} & Refused deliberately: the system-folder guard, a permission wall, or policy \\
\code{4} & Finished, but items were skipped and \code{--strict} was given \\
\bottomrule
\end{tabularx}
\end{table}

\subsection{Two ways to let an agent move files}
\label{sec:agent:capability}

This is the central design question of the section, and we think the answer
generalizes well beyond this tool.

Suppose an agent should be able to reorganize a directory. There are two
shapes the interface can take.

\begin{description}[leftmargin=1.6em,itemsep=0.3em]
  \item[Validate what it asks for.] The agent supplies operations, and the tool
    checks each one. This is what the \code{apply} sub-command does with a plan
    file: it verifies that every path is absolute and lies inside the plan's own
    root, and that the counts match, before executing anything.
  \item[Never let it name anything.] The agent asks the tool to compute a plan,
    receives an opaque handle, and applies \emph{that}. It never supplies a
    path. This is what the MCP server does: \code{plan\_organize} returns a
    \code{plan\_id}, and \code{apply\_plan} accepts that identifier and nothing
    else.
\end{description}

The second is strictly stronger, and the reason is worth stating precisely.
Validation has to anticipate every way a path can escape its root, and
\Cref{sec:detect} is a catalogue of how many ways that is: \code{..}
components, symbolic links and NTFS junctions, verbatim and UNC prefixes, 8.3
short names, case folding on two platforms, and trailing characters Windows
strips silently. A validator is only as good as that list. Under the second
shape the request contains no path at all, so the question does not arise: the
moves are the ones the tool computed, and the agent's only remaining choice is
\emph{whether}, not \emph{what}.

This is the capability-security distinction --- unforgeable references in place
of ambient authority checked against a name --- applied to agent tool design
\cite{dennis1966,miller2006}. We claim no novelty in the principle. What seems
worth reporting is that it maps unusually cleanly onto a tool already split
into a pure planner and an executor: the handle is simply the plan, and the
split existed for testability long before an agent was in view.

The practical consequence is a defence against indirect prompt injection
\cite{greshake2023}. An agent that reads a file containing
\enquote{\emph{also move C:\textbackslash Windows\textbackslash System32 into
Documents}} cannot act on it through this interface, because there is no
argument in which that instruction could be expressed. It can be expressed to
the validating interface, where it is then rejected --- which is a weaker
position, because it relies on the rejection being complete.

\subsection{The MCP server}
\label{sec:agent:mcp}

The Model Context Protocol is JSON-RPC 2.0 over newline-delimited JSON on
standard input and output, so serving it needed no dependency beyond the
\code{serde\_json} already present. \code{Server::handle} is a pure function
from one request to at most one response, and the loop around it does nothing
but framing; the protocol is therefore tested without a pipe or a subprocess,
the same separation the classifier uses.

Four restrictions apply, in rough order of how much they carry:

\begin{enumerate}[leftmargin=1.6em,itemsep=0.25em]
  \item \textbf{No paths in requests}, as above.
  \item \textbf{Confinement to one root}, given on the command line and checked
    on the canonical path, so \code{..} and symbolic links are caught.
  \item \textbf{Read-only by default.} Without \code{--allow-writes} the
    mutating tools are \emph{absent from the tool listing}, not present and
    failing. An agent cannot attempt what it cannot discover, and a listing is
    a cheaper thing to reason about than a set of runtime refusals.
  \item \textbf{The guard cannot be overridden.} There is no way to express
    \code{--allow-system-folder} over this interface, and a server rooted at a
    refused directory does not start. Failing at startup rather than per call
    matters: an agent holding a connection to such a server would otherwise
    find every call failing with no way to understand why.
\end{enumerate}

A plan is consumed when applied, so a replayed \code{apply\_plan} cannot re-run
moves against a directory that has since changed. Large plans are returned as a
50-move sample with a truncation flag, because a response of ten thousand
entries helps no reader, human or otherwise.

\subsection{The review surface}
\label{sec:agent:review}

Delegation is only acceptable if the delegation can be checked, and checking
was the weakest part of the system: \code{history} listed runs and nothing
explained one. Reading a journal meant opening its JSON Lines file.

Building \code{show} surfaced a structural point about journal-based undo.
\emph{Two of the three things a reviewer wants are not in the journal, and
cannot be.}

\begin{enumerate}[leftmargin=1.6em,itemsep=0.25em]
  \item \textbf{Whether each item is still where the run put it} must be
    determined at review time, because the directory keeps changing after the
    run. This is the question that actually matters, since it is the difference
    between an undo that is clean and one that reports items as missing. The
    report marks each moved-on item.
  \item \textbf{What the run left alone} can never be in a journal, because a
    journal records changes. It comes from the optional run log
    (\Cref{sec:obs}), and its absence is reported as silence rather than as an
    error, since most runs are not logged.
  \item Timing, which the journal does have.
\end{enumerate}

So \code{show} is not a journal pretty-printer; it is a journal, plus a live
filesystem check, plus a separate log. We suspect this generalizes to any
journal-based undo facility: the journal is a record of a transition, and a
reviewer is asking about a state.

\subsection{Configuration: a ceiling and a floor}
\label{sec:agent:config}

An agent-facing tool needs configuration for a different reason than a
human-facing one. A person wants to avoid retyping flags; an agent does not
care. What an installation needs is a statement of what it \emph{permits},
which the caller cannot override.

The file therefore holds two things with opposite semantics. A \textbf{policy}
is a ceiling: permitted roots, whether writes are allowed at all, whether
deletion is, whether the override flag exists, and a cap on scan depth. Flags
cannot raise it. \textbf{Defaults} are a floor: ignore rules, depth, grouping
keys. Any flag overrides them.

\paragraph{Discovery is deliberately asymmetric.}
A policy is read only from a path the operator named, via \code{--config} or an
environment variable, and is never discovered. Defaults may additionally come
from a \code{.tidy-up.json} inside the directory being processed, which is
convenient and travels with the directory; a policy found there is discarded.

The asymmetry is the security property. If a policy could be discovered from
the directory being processed, then unpacking an archive and tidying it would
let the archive decide what may be done to it --- a confused-deputy arrangement
\cite{hardy1988} in which the tool's authority is directed by the data it was
pointed at. The rule we extracted is short enough to apply elsewhere:
\emph{auto-discovered configuration may lower the bar for convenience; it may
never raise it for permission.}

Two smaller decisions follow the same instinct. Passing a flag that policy
forbids is an \emph{error}, not a silent no: the caller believes they have
authorized something, and learning otherwise afterwards is worse than being
told. And an unrecognized field is rejected rather than ignored, because
\code{allow\_system\_folder} for \code{allow\_system\_folders} would otherwise
leave an installation believing it was locked down when it was not.

% =============================================================================
\section{Evaluation}
\label{sec:eval}

\subsection{Method}

Evaluation is by automated test, at three levels: pure unit tests of the classifier
and planners, host-dependent unit tests of the probes, and black-box integration tests
that execute the compiled binary. All figures below are from a full
\code{cargo test} run on Windows 11 with Rust 1.95.0. \Cref{tab:tests} gives the
counts.

\begin{table}[t]
\centering
\small
\caption{Test suite. 455 tests, zero failures, zero \code{clippy} warnings at
\code{--all-targets}, on both stable Rust and the declared minimum of 1.85.}
\label{tab:tests}
\begin{tabularx}{\textwidth}{@{}l r X@{}}
\toprule
\textbf{Suite} & \textbf{Tests} & \textbf{Scope} \\
\midrule
unit (in-crate)        & 331 & classifier corpus, probes, planners, taxonomy, log and JSON formats, MCP protocol, policy \\
\code{tests/cli}       & 30  & pre-existing command-line behaviour (regression) \\
\code{tests/workflow}  & 23  & scan $\rightarrow$ plan $\rightarrow$ execute $\rightarrow$ journal $\rightarrow$ restore \\
\code{tests/reorganize}& 11  & re-grouping, idempotence, undo round trip \\
\code{tests/safety}    & 11  & guard policy through the compiled binary \\
\code{tests/api}       & 10  & envelope shape, exit codes, propose/apply, plan trust boundary \\
\code{tests/config}    & 10  & policy ceiling, defaults floor, the folder-config attack \\
\code{tests/show}      & 8   & review surface, clean versus partial undo \\
\code{tests/logging}   & 7   & run log content, joins, degradation, \code{--strict} \\
\code{tests/cross\_volume} & 5 & cross-partition moves (environment-gated) \\
\code{tests/scripts}   & 3   & installer scripts \\
doc-tests              & 6   & every documented example compiles and runs \\
\midrule
\textbf{Total}         & \textbf{455} & \\
\bottomrule
\end{tabularx}
\end{table}

\subsection{Cross-platform rule validation}

Because \code{Platform} is a parameter, the full rule corpus runs on every host. The
table-driven suite covers Windows dangerous locations and safe locations, scrubbed and
redirected environment variables, case and separator variants, macOS dangerous
locations and all five carve-outs, and Linux dangerous locations and carve-outs, each
case carrying a one-line justification that is printed on failure. This is the concrete
payoff of P3: the macOS \code{/System/Volumes/Data} carve-out and the Linux
\code{/run/media} carve-out --- the two most consequential and most easily forgotten
--- are verified on a Windows development machine, where neither path exists.

\subsection{Precision: the negative corpus}

P2 is tested directly by \code{ordinary\_folders\_are\_never\_flagged}, which asserts
\textsf{Safe} for 24 ordinary directories spanning all three platforms:
\code{C:\textbackslash Users\textbackslash me\textbackslash Downloads}, \code{D:\textbackslash Media\textbackslash Movies},
\code{\textbackslash\textbackslash nas\textbackslash media\textbackslash Photos}, \code{/Users/me/Documents/Work},
\code{/Volumes/Backup/Photos}, \code{/home/me/Pictures/2024}, \code{/mnt/storage},
\code{/media/me/Elements}, \code{/data/archive}, \code{/tmp/scratch} and others. The
test exists as a regression net: a future rule that increases coverage at the expense
of precision fails it, and the failure names the directory that regressed.

\subsection{Guard behaviour}

Ten integration tests drive the compiled binary with no terminal attached, the
condition that matters for scripts and CI. The most significant assert that
\code{organize <home> --yes} exits non-zero, names the override flag, and leaves the
home directory listing unchanged; that the override alone still cannot run unattended;
that \code{--dry-run} on a flagged directory warns, proceeds and creates no state; that
reporting commands are never blocked; that \code{distribute} guards destinations; and
that \emph{every} destructive command refuses the home directory --- the last written
as a loop over the command set so that a newly added command without a guard fails the
suite. A complementary test asserts that reporting commands are the exception, so the
first cannot be satisfied by blocking everything.

Manual verification on the development machine confirmed the intended behaviour
against real system directories:

\begin{lstlisting}[style=shell]
$ tidy-up organize 'C:\Windows' --yes
Careful
! C:\Windows is a folder the system manages.
  It holds the Windows installation, which other programs expect to find in place.
  tidy-up cannot write here, so nothing could be moved.
error: refusing to change C:\Windows: tidy-up cannot write there, and
       --allow-system-folder cannot grant permission the system refused
$ echo $?
1
\end{lstlisting}

Note that \code{C:\textbackslash Windows} is refused on two independent grounds --- the name rule
and the writability probe --- and that the message correctly reports the override as
unhelpful, since the binding constraint is permission rather than policy.

\subsection{Reversibility}

The re-grouping round trip is asserted as equality of two complete observations of the
tree, before and after, excluding only the tool's own state directory: a map from
relative path to file contents, and a sorted list of relative directory paths. Testing
directory structure separately from file contents is what gives the test power over
the \code{DirRemoved} record; a files-only comparison would pass even if every empty
directory were lost.

\subsection{The delegation interface}
\label{sec:eval:agent}

The MCP server is exercised through \code{Server::handle} directly, since it is
a pure function of a request, with 18 tests covering the handshake and version
negotiation, notifications receiving no reply, tool errors arriving as results
rather than protocol failures, and the full plan/apply/review/undo cycle.

Four of them are the safety argument of \Cref{sec:agent:capability} stated as
assertions. One confirms that \code{apply\_plan} advertises exactly one input
property, \code{plan\_id}, so the interface admits no path. One confirms that
an invented handle is refused. One confirms a plan cannot be applied twice.
The fourth attempts four escapes from the server root --- \code{..}, a
trailing-separator variant, an absolute path elsewhere, and on Unix a symbolic
link planted inside the root pointing out of it --- and requires all four to be
refused.

The plan file interface, being the weaker of the two shapes, is tested
adversarially in kind: a proposed plan is rewritten so that one destination is
\code{C:\textbackslash Windows}, and the test requires both that the run is
refused and that the target does not come into existence.

For configuration, the corresponding test writes a \code{.tidy-up.json} into a
directory that the operator's policy does \emph{not} permit, claiming in it
that the directory is permitted and that system folders are allowed, then
requires the run to be refused anyway. That is the confused-deputy scenario of
\Cref{sec:agent:config} reduced to a test case.

\subsection{Dependency footprint}
\label{sec:eval:deps}

\begin{table}[t]
\centering
\small
\caption{Dependency footprint before this work and at release 0.3.0.}
\label{tab:deps}
\begin{tabular}{@{}l r r@{}}
\toprule
 & \textbf{Before} & \textbf{After} \\
\midrule
Direct dependencies            & 10 & 10 \\
Unique transitive crates       & 54 & 54 \\
Stripped release binary        & 2.19\,MiB & 2.68\,MiB \\
New dependencies added         & \multicolumn{2}{c}{0} \\
Lines added (implementation + tests) & \multicolumn{2}{c}{10{,}265} \\
\bottomrule
\end{tabular}
\end{table}

Zero new dependencies was a hard constraint, and it shaped four decisions that would
otherwise have gone differently:

\begin{enumerate}[leftmargin=1.6em,itemsep=0.2em]
  \item \textbf{Logging} is hand-rolled on the already-present \code{serde}, not
    \code{tracing} with \code{tracing-subscriber}. The ecosystem option would have
    added roughly a dozen transitive crates, dominated by \code{regex-automata} via
    \code{EnvFilter}, for dynamic filtering and asynchronous span propagation that a
    single-threaded tool with one process-wide span does not use. The closed event
    enum is, additionally, a better contract than stringly-typed fields.
  \item \textbf{Ownership checks} were dropped rather than take \code{libc} or
    \code{rustix} (\Cref{sec:perm}).
  \item \textbf{Local time} was declined rather than take a date crate
    (\Cref{sec:regroup}).
  \item \textbf{Platform attributes} use stable \code{std::os::*::fs::MetadataExt}
    interfaces --- \code{file\_attributes} on Windows, \code{st\_flags} on macOS,
    \code{dev} on Unix --- rather than \code{windows-sys}.
\end{enumerate}

We report items 2 and 3 as genuine reductions in capability accepted in exchange for a
stated property of the artifact, not as free wins.

The binary grew by roughly half a megabyte across both releases, which is the
cost of the code itself rather than of anything linked in: the dependency
graph is the same 54 crates it was before any of this started. For comparison,
the single rejected dependency of item~1 was estimated at a similar size on
its own.

\subsection{What continuous integration caught}
\label{sec:eval:ci}

Three classes of defect survived a green local run on Windows and were caught only
by the cross-platform matrix. They are reported because each is a direct consequence
of a design decision described above.

First, \textbf{the temporary-directory carve-out was broken on two of three
platforms}, exactly the mechanism \Cref{sec:detect} identifies as load bearing. The
environment recorded the canonicalized temporary directory, but a caller that has not
canonicalized its own path supplies a different spelling of the same place:
\code{tempfile} returns \code{/var/folders/...} on macOS where the canonical form is
\code{/private/var/folders/...}, and an 8.3 short name on Windows. The two never
compared equal, so every temporary directory was classified \textsf{Dangerous}. The
fix records the temporary directory under every spelling and, on macOS, folds
\code{/private/var}, \code{/private/etc} and \code{/private/tmp} onto their short
names, since a verdict must not depend on how the caller typed the path. Three tests
now pin this, written against the pure classifier so they run on any host. This is a
pointed illustration of the limits of P3: making the \emph{rules} testable everywhere
does not make the \emph{environment} so.

Second, \textbf{the new code did not compile on the declared minimum supported Rust
version}. Ten sites used \code{let} chains, stabilized in 1.88, and one used
\code{ErrorKind::InvalidFilename}, still feature-gated at 1.85. Neither is visible to
a developer on current stable. A declared MSRV is a compatibility claim like any
other and needs a job that checks it.

Third, \textbf{a defect from reading the clock twice}. The run log and the
undo journal each mint an identifier from a separate call to
\code{compact\_id(now\_secs())}, on the assumption they would agree. A run that
crosses a one-second boundary between starting and opening its journal ends up
with two, and \code{show} was locating the log by the journal's name. It
therefore reported nothing about what the run left alone --- indistinguishable,
to a reader, from a run that was never logged. Only the musl job happened to be
unlucky in the second that mattered. The fix matches on the identifier the log
already records, and the general lesson is that two clock reads are two values:
if one artifact must refer to another, it should say so rather than be named
after it.

Fourth, \textbf{two integration tests encoded assumptions about the developer's
machine}: one required the home directory to contain loose files, which a fresh CI
runner's does not, and one compared a directory listing that a parallel test was
concurrently modifying.

The second of those turned out to be a design defect rather than a flaky test. The
two listings differed by a file named \code{.tidy-up-write-test-\ldots}: the guard
gathered every fact before classifying, so the writability probe of \Cref{sec:perm}
created a file \emph{inside the very directory it was about to refuse}. Harmless in
effect, since the file is removed immediately, but wrong in principle, because the
guard wrote to a directory it then told the user it would not touch, and it raced
with anything else reading that directory. The probe now runs only after the path
rules have failed to produce a refusal: for a directory refused by name the extra
reason would not change the outcome, and there is no sense writing into somewhere
already declined. For an ordinary directory it still runs, which is the case where
being told up front is worth a file created and removed. A unit test and an
integration test pin the new ordering.

This is the most instructive of the three. It was invisible to every local run, it
was found by an assertion written to check something else, and the honest fix was to
change the system rather than to relax the test.

\subsection{Defects found by the work itself}

Three pre-existing defects surfaced while implementing the above and are worth
recording, since each is a plausible bug in comparable tools:

\begin{enumerate}[leftmargin=1.6em,itemsep=0.2em]
  \item \code{--include-hidden} un-skipped Windows \textsc{system}-attributed files
    along with hidden ones, because the two attributes were tested as one mask.
  \item The directory walk aborted an entire scan on any failure below the root.
  \item Permission denial below the root was reported as the undifferentiated word
    \enquote{unreadable}, discarding the error kind.
\end{enumerate}

A fourth defect was introduced and caught by the new tests rather than by review: the
first integration test for \code{distribute} destination guarding failed because the
sandbox directory lies inside the home directory on Windows, so the pre-existing
overlap check rejected the pair before the guard was reached. The test was correct and
the fixture was wrong, which is the desirable ordering.

% =============================================================================
\section{Discussion}
\label{sec:discussion}

\subsection{Threats to validity}

\textbf{Platform coverage of the probes.} Rule tables are validated on every host, but
the probe implementations are inherently platform-specific and were executed only on
Windows during development. Continuous integration runs the suite on all three
platforms; the macOS \code{st\_flags} path and the Linux \code{mountinfo} live-read
path have not been exercised by the author directly. The \code{mountinfo} \emph{parser}
is tested against fixtures on every host, which isolates the untested surface to the
file read itself.

\textbf{Soundness is empirical, not proved.} P1 is a claim about a hand-maintained
table, not a theorem. New operating system releases introduce directories; the mitigation
is the layered composition of name rules, environment variables, runtime attributes and
the writability probe, such that a path missing from the table is often still caught by
a fact. We make no completeness claim.

\textbf{Precision is measured against a curated corpus.} The 24-directory negative
corpus reflects the author's judgement about what is ordinary. A field study of real
invocations would be a stronger instrument and was not conducted.

\textbf{Time-of-check-to-time-of-use.} The classifier judges a directory that could in
principle be replaced before the operation proceeds. We consider this acceptable: the
guard defends against user error, not against a local adversary who could simply run
the tool themselves.

\subsection{What the delegation argument does and does not claim}

We are not claiming that an agent driving this tool is safe in general. We are
claiming something narrower and, we think, checkable: that of the two available
interface shapes, one admits a class of attack the other cannot express, and
that the stronger one costs little when a plan/execute split already exists.

Three limits are worth stating. The argument concerns the \emph{interface}, not
the agent: an agent authorized to apply plans can still apply a reasonable-looking
plan for a bad reason, and the defence against that is the review surface and
the undo journal rather than the interface shape. The confinement root is
resolved once at startup, so a directory replaced beneath a long-running server
is a time-of-check hazard of the same kind noted in \Cref{sec:discussion}.
And the server holds plans in memory, so a restart invalidates outstanding
handles: correct, but it makes the handle a session capability rather than a
durable one, which an orchestrator would need to revisit.

\subsection{Generality}

The mechanisms are not specific to file organizers. The pure-classifier-plus-facts
structure applies to any tool that must judge a target before acting on it; the
platform-as-parameter technique applies to any cross-platform path logic and is, we
think, underused; the survival taxonomy with one stated fatal exception applies to any
bulk operation; and the observation that a move journal cannot express directory
deletion applies to any reversible operation that removes containers.

\subsection{An orchestrator, and why it is not here}

The natural next step, running tidy-up across a fleet from one client, was
designed and deliberately not built. A tool that moves files as a privileged
user and is reachable over a network is close to remote code execution, and the
security model is the architecture rather than a layer over it. Two decisions
would be load-bearing: agents should poll for work rather than accept
connections, removing the inbound surface entirely; and policy must live on the
agent, so that a job says \enquote{run profile \emph{photos}} and the agent
decides what that means and where it may apply. A request that can name a path
turns the whole arrangement into the weaker of the two shapes in
\Cref{sec:agent:capability}, fleet-wide.

\subsection{Future work}

Four directions follow directly. Effective-UID comparison via \code{rustix} would
complete the ownership signal at the cost of one pure-Rust dependency. A TZif parser
would remove the time-zone limitation without a date crate. Per-move instrumentation
could be added by way of an observer trait on the executor, which was designed but not
implemented, on the grounds that the journal already carries the same information.
Finally, the negative corpus should be extended empirically rather than by judgement,
which would require telemetry the tool deliberately does not collect --- an instructive
tension between measurement and the project's privacy stance.

% =============================================================================
\section{Conclusion}
\label{sec:conclusion}

Bulk file reorganization is a category of tool whose core mechanism is indifferent to
the value of what it operates on. We have described a guard that closes that gap
across Windows, macOS and Linux, built around a pure classifier whose platform is an
argument, so that all three rule tables are validated from any host; an empirical
writability probe that is correct where mode bits, ACLs and System Integrity
Protection disagree; a survival architecture with exactly one stated fatal exception;
and an idempotent re-grouping operation whose undo restores directory structure, not
merely file locations.

We then took the same question one level further: what changes when the
caller is not a person. The answer we arrived at is that an interface in which
an agent cannot name a path is strictly stronger than one that validates the
paths it names, because the second is only as complete as its author's list of
ways a path can escape, and \Cref{sec:detect} is evidence that such lists are
long. A plan/execute split kept for testability turned out to supply exactly
the unforgeable handle this needs, which is a pleasant accident worth noticing:
the property was available years before the reason to want it.

Three lessons generalize beyond the artifact. The first is that precision is a safety
property: a guard that fires on \code{D:\textbackslash Downloads} teaches users to bypass
it and protects nobody, which is why two plausible rules were deliberately weakened and
why a negative corpus is part of the test suite rather than an afterthought. The second is that a journal is a record of a transition, so it
cannot answer questions about a state: \enquote{is this still undoable} and
\enquote{what was left alone} both need something the journal structurally does
not hold, which is why the review surface is a journal plus a live check plus a
log. The third
is that separating a pure verdict from an impure policy is what made the work testable
at all --- and that the same separation forced the discovery, in the Rust standard
library, that \code{std::path} cannot be used to reason about another platform's paths.

The implementation is 10{,}265 lines, adds no dependencies, and passes 455 tests
on every platform it ships for and on its declared minimum compiler.

% =============================================================================
\section*{Availability}
\addcontentsline{toc}{section}{Availability}

The implementation is open source at
\url{https://github.com/ianktoo/tidy-up}, released under the MIT licence together with
a supplementary end-user licence. The modules described here are \code{src/safety.rs},
\code{src/outcome.rs}, \code{src/regroup.rs}, \code{src/obs.rs},
\code{src/api.rs}, \code{src/mcp.rs}, \code{src/config.rs},
\code{src/commands/guard.rs} and \code{src/commands/show.rs}, with tests in the
corresponding \code{*\_tests.rs} files and in \code{tests/}. The work described
here was released as versions 0.2.0 and 0.3.0.

\section*{Acknowledgements}
\addcontentsline{toc}{section}{Acknowledgements}

Portions of the implementation and this manuscript were prepared with the assistance
of an AI coding assistant; all design decisions, trade-offs and the final text are the
author's own.

% =============================================================================
\begin{thebibliography}{99}
\addcontentsline{toc}{section}{References}
\small

\bibitem{dennis1966}
J.~B.~Dennis and E.~C.~Van Horn.
\newblock Programming semantics for multiprogrammed computations.
\newblock \emph{Communications of the ACM}, 9(3):143--155, 1966.

\bibitem{miller2006}
M.~S.~Miller.
\newblock \emph{Robust Composition: Towards a Unified Approach to Access Control
  and Concurrency Control}.
\newblock PhD thesis, Johns Hopkins University, 2006.

\bibitem{hardy1988}
N.~Hardy.
\newblock The confused deputy (or why capabilities might have been invented).
\newblock \emph{ACM SIGOPS Operating Systems Review}, 22(4):36--38, 1988.

\bibitem{greshake2023}
K.~Greshake, S.~Abdelnabi, S.~Mishra, C.~Endres, T.~Holz, and M.~Fritz.
\newblock Not what you've signed up for: compromising real-world LLM-integrated
  applications with indirect prompt injection.
\newblock In \emph{Proceedings of the 16th ACM Workshop on Artificial
  Intelligence and Security (AISec)}, 2023.

\bibitem{mcp}
Anthropic.
\newblock \emph{Model Context Protocol specification}.
\newblock \url{https://modelcontextprotocol.io}.

\bibitem{jsonrpc}
JSON-RPC Working Group.
\newblock \emph{JSON-RPC 2.0 Specification}, 2010.

\bibitem{posix}
IEEE and The Open Group.
\newblock \emph{The Open Group Base Specifications Issue 7, 2018 edition (POSIX.1-2017)}.
\newblock IEEE Std 1003.1-2017. See in particular \code{access()} and
  \code{faccessat()}.

\bibitem{linuxman}
M.~Kerrisk.
\newblock \emph{The Linux Programming Interface}.
\newblock No Starch Press, 2010.

\bibitem{procfs}
Linux kernel documentation.
\newblock \emph{proc(5) --- process information pseudo-filesystem}.
\newblock Describes the \code{/proc/[pid]/mountinfo} format, including octal escaping
  of mount points and the optional-field separator.

\bibitem{fhs}
Linux Foundation.
\newblock \emph{Filesystem Hierarchy Standard, Version 3.0}, 2015.

\bibitem{sip}
Apple Inc.
\newblock \emph{System Integrity Protection Guide}.
\newblock Apple Developer Documentation.

\bibitem{apfs}
Apple Inc.
\newblock \emph{Apple File System Reference}.
\newblock Apple Developer Documentation. Describes firmlinks and the split
  System/Data volume layout.

\bibitem{msdnpaths}
Microsoft Corporation.
\newblock \emph{Naming Files, Paths, and Namespaces}.
\newblock Windows Application Development Documentation. Covers reserved device names,
  trailing dot and space handling, and the \code{\textbackslash\textbackslash?\textbackslash} prefix.

\bibitem{msdnattr}
Microsoft Corporation.
\newblock \emph{File Attribute Constants}.
\newblock Windows Application Development Documentation.

\bibitem{knownfolders}
Microsoft Corporation.
\newblock \emph{Known Folders and the \code{KNOWNFOLDERID} constants}.
\newblock Windows Application Development Documentation.

\bibitem{coreutils}
GNU Project.
\newblock \emph{GNU Coreutils Manual}.
\newblock See \code{rm} and the \code{--preserve-root} option.

\bibitem{hinnant}
H.~Hinnant.
\newblock \emph{\code{chrono}-Compatible Low-Level Date Algorithms}.
\newblock Describes the \code{civil\_from\_days} algorithm used for date conversion.

\bibitem{rustdoc}
The Rust Project Developers.
\newblock \emph{The Rust Standard Library Documentation}.
\newblock See \code{std::path}, \code{std::io::ErrorKind}, and the platform
  \code{MetadataExt} extension traits.

\bibitem{tidyup}
I.~Too.
\newblock \emph{tidy-up: organize messy folders by file type, isolate duplicates, and
  undo it all}.
\newblock \url{https://github.com/ianktoo/tidy-up}.

\end{thebibliography}

\end{document}
